\section{Experimental Setup}
\label{sec:experiments}

\subsection{Datasets and Tasks}

We evaluated the RCD-hosted \texttt{qwen3.5-9b} model on the anatomical
recognition and diagnosis tasks derived from LMOD~\cite{qin2025lmod}.  The
evaluation included 7,245 images: 3,859 optical coherence tomography (OCT)
images from OIMHS and 3,386 color fundus photographs from REFUGE
($n=1{,}200$), ORIGA ($n=650$), G1020 ($n=1{,}020$), and IDRiD ($n=516$).
Across these images, the anatomical recognition task contained 15,938
reference regions.  The model was shown the bounding-box-annotated image and
asked to return one anatomical type for each region identifier using the
format \texttt{Region ID: <ID>; Type: <TYPE>}.

For diagnosis, we evaluated glaucoma detection on all 2,470 fundus images
with explicit glaucoma labels and macular-hole staging on all 3,859 OIMHS
images.  The glaucoma set contained 544 glaucoma and 1,926 non-glaucoma
images.  Its source distribution was 800 REFUGE, 650 ORIGA, and 1,020 G1020
images.  The macular-hole reference distribution was 19, 610, 1,043, and
2,187 images for stages 1 through 4, respectively.  We evaluated every
available reference-labeled example without class balancing or resampling.
Consequently, the reported glaucoma accuracy should be interpreted together
with the class distribution and majority-class baseline.

\subsection{Inference Configuration}

We accessed \texttt{qwen3.5-9b} through the Clemson RCD OpenAI-compatible
multimodal endpoint.  Local images were converted to RGB JPEG at quality 95
and resized, when necessary, such that the longest side did not exceed 2,048
pixels.  We used deterministic decoding with temperature 0, requested that
model thinking be disabled, and allowed at most 1,024 completion tokens.  The
request timeout was 600 seconds with at most five client retries.  Predictions
were checkpointed after each request.  Clean images were used for diagnosis,
whereas annotated images were used for anatomical recognition.

The recognition prompt listed the allowed anatomical categories for the
corresponding modality.  For diagnosis, glaucoma responses were requested in
the form \texttt{GLAUCOMA / NON-GLAUCOMA; Explanation: <TEXT>}, and
macular-hole responses were requested as
\texttt{Stage: <INTEGER>; Justification: <TEXT>}.  We reparsed the raw model
responses during evaluation rather than relying on intermediate prediction
fields.  A macular-hole response was valid only when the explicit
\texttt{Stage:} field contained an integer from 1 to 4.  This prevented a
number appearing later in a free-text explanation from being mistaken for the
predicted stage.

\subsection{Evaluation Metrics}

For anatomical recognition, let $C$ be the number of predicted region
identifiers whose normalized anatomical type matches the LMOD reference, let
$N_P$ be the number of unique predicted region identifiers, and let $N_T$ be
the number of reference region identifiers.  We computed micro-averaged
precision, recall, and F1 as
\begin{equation}
    \mathrm{Precision}=\frac{C}{N_P},\qquad
    \mathrm{Recall}=\frac{C}{N_T},\qquad
    \mathrm{F1}=\frac{2\,\mathrm{Precision}\,\mathrm{Recall}}
    {\mathrm{Precision}+\mathrm{Recall}}.
\end{equation}
Following LMOD~\cite{qin2025lmod}, hallucination resistance for image $i$ was
computed from the predicted identifier set $\mathcal{P}_i$ and reference
identifier set $\mathcal{T}_i$:
\begin{equation}
    \mathrm{HR}_i = 1-
    \frac{\left|\left\{r\in\mathcal{P}_i:r\notin\mathcal{T}_i\right\}\right|}
    {|\mathcal{P}_i|}.
\end{equation}
We report the macro mean of $\mathrm{HR}_i$ over parseable responses.  HR is
undefined for an empty predicted set; these cases were instead captured by
the invalid-response rate.

For glaucoma detection and macular-hole staging, accuracy was the number of
correct strict-format predictions divided by the total number of
reference-labeled examples.  Missing and invalid responses therefore counted
as incorrect.  For each task, invalid-response rate (IR) was
\begin{equation}
    \mathrm{IR}=\frac{N_{\mathrm{missing}}+N_{\mathrm{unparseable}}}
    {N_{\mathrm{reference}}}.
\end{equation}
All metrics were computed directly from the saved raw predictions and the
\texttt{information.json} reference annotations in the extracted LMOD data.

\section{Results}
\label{sec:results}

\subsection{Overall Performance}

Table~\ref{tab:results} summarizes performance across the two tasks.
Anatomical recognition achieved 64.43\% precision, 63.03\% recall, and a
63.72\% F1 score.  Of 15,591 parsed region predictions, 10,045 matched both
the reference identifier and anatomical type.  The model produced only 53
invented region identifiers, yielding an HR of 99.76\%.  Thus, the main
recognition errors arose from omitted or incorrectly typed reference regions
rather than from inventing new identifiers.

The recognition invalid-response rate was 2.40\% (174 of 7,245 images).  This
total consisted of 131 unparseable outputs and 43 missing prediction records.
The current recognition checkpoint therefore contains 7,202 of 7,245 expected
records; missing records are included as invalid and contribute no correct
regions in the reported metrics.

For glaucoma detection, the model obtained 77.17\% accuracy (1,906 of 2,470)
with a 0.20\% invalid-response rate.  Accuracy among the 2,465 valid responses
was 77.32\%.  However, the majority-class baseline obtained by always
predicting non-glaucoma is 77.98\% on this imbalanced evaluation set.  The
model predicted non-glaucoma for 2,087 valid images and glaucoma for only 378,
indicating a strong bias toward the majority class.  Therefore, the apparently
high glaucoma accuracy does not establish discrimination beyond a trivial
majority classifier.

Macular-hole staging was substantially more difficult.  Overall accuracy was
15.39\% (594 of 3,859), and 16.30\% of responses were invalid.  Even after
excluding the 629 invalid responses, accuracy was only 18.39\% (594 of 3,230),
below the 25\% uniform-random reference for a four-class task.  Among valid
outputs, the model predicted stages 1, 2, 3, and 4 for 969, 715, 1,429, and
117 images, respectively.  This distribution contrasts sharply with the
reference set, which is dominated by stage 4 ($n=2{,}187$), and helps explain
the low staging accuracy.

\begin{table*}[t]
    \centering
    \caption{Current performance of RCD-hosted \texttt{qwen3.5-9b} on the
    evaluated LMOD tasks. All values are percentages. Invalid or missing
    diagnosis responses count as incorrect in accuracy. Recognition HR is
    averaged over parseable responses.}
    \label{tab:results}
    \small
    \begin{tabular}{lrrrrrrrrr}
        \toprule
        & \multicolumn{5}{c}{Anatomical recognition}
        & \multicolumn{2}{c}{Glaucoma}
        & \multicolumn{2}{c}{MH staging} \\
        \cmidrule(lr){2-6}\cmidrule(lr){7-8}\cmidrule(lr){9-10}
        Model & Prec. & Rec. & F1 & HR & IR & Acc. & IR & Acc. & IR \\
        \midrule
        \texttt{qwen3.5-9b}
        & 64.43 & 63.03 & 63.72 & 99.76 & 2.40
        & 77.17 & 0.20 & 15.39 & 16.30 \\
        \bottomrule
    \end{tabular}
\end{table*}

\subsection{Recognition Performance by Modality}

Recognition performance varied substantially by modality
(Table~\ref{tab:recognition_by_modality}).  On the pooled CFP group, comprising
REFUGE, ORIGA, G1020, and IDRiD, the model achieved 97.37\% F1, compared with
40.02\% on OIMHS OCT.  The OIMHS subset accounted for 133 of the 174 invalid
recognition responses and 47 of the 53 invented region identifiers.  Thus, the
overall recognition score masks a pronounced modality gap: labeled optic-disc,
optic-cup, and foveal regions in fundus photographs were recognized reliably,
whereas retinal, choroidal, intraretinal-cyst, and macular-hole regions in OCT
remained challenging.

\begin{table}[t]
    \centering
    \caption{Anatomical recognition performance by imaging modality. CFP
    pools REFUGE, ORIGA, G1020, and IDRiD. All values are percentages.}
    \label{tab:recognition_by_modality}
    \small
    \begin{tabular}{lrrrrr}
        \toprule
        Modality & Prec. & Rec. & F1 & HR & IR \\
        \midrule
        OCT (OIMHS) & 40.64 & 39.42 & 40.02 & 99.58 & 3.45 \\
        CFP (others) & 97.87 & 96.87 & 97.37 & 99.96 & 1.21 \\
        \bottomrule
    \end{tabular}
\end{table}

\subsection{Summary of Findings}

Three findings emerge from the current evaluation.  First, the near-perfect
HR indicates that \texttt{qwen3.5-9b} generally respected the set of region
identifiers presented in the prompt, although this did not guarantee correct
anatomical classification.  Second, recognition performance was highly
modality dependent, with a 57.35-point F1 gap between OCT and pooled CFP.
Third, diagnosis remained unreliable: glaucoma
accuracy did not exceed the majority baseline, and macular-hole staging was
below the uniform-random reference even when invalid responses were excluded.
Together, these findings suggest that output-format compliance and resistance
to identifier hallucination are necessary but insufficient for clinically
meaningful ophthalmic reasoning.
